\section*{Capítulo \ref{ch:b2:ligand-field} --- Teoria do campo ligante e cor}

\begin{solution}{exo:b2:ligand-field:1}
$d^3$: $t_{2g}^3$, EECC $-1.2\Delta_o$. $d^8$: $t_{2g}^6e_g^2$, EECC $-2.4 + 1.2 = -1.2\Delta_o$.
\end{solution}

\begin{solution}{exo:b2:ligand-field:2}
\ce{[Fe(H2O)6]^2+}, $d^6$ de \omterm{def:b2:ligand-field:spin}{spin alto} $t_{2g}^4e_g^2$: quatro. \ce{[Fe(CN)6]^4-}, \omterm{def:b2:ligand-field:spin}{spin baixo}
$t_{2g}^6$: nenhum (\omterm{def:b2:diatomic-mos:magnetism}{diamagnético}).
\end{solution}

\begin{solution}{exo:b2:ligand-field:3}
600 nm é luz laranja; o complexo parece azul (azul-esverdeado).
\end{solution}

\begin{solution}{exo:b2:ligand-field:4}
$\ce{Cl-} < \ce{H2O} < \ce{NH3} < \ce{CN-}$.
\end{solution}

\begin{solution}{exo:b2:ligand-field:5}
Água: $\Delta_o < P$, \omterm{def:b2:ligand-field:spin}{spin alto}, $t_{2g}^3e_g^2$, cinco elétrons desemparelhados. Cianeto: $\Delta_o > P$,
\omterm{def:b2:ligand-field:spin}{spin baixo}, $t_{2g}^5$, um elétron desemparelhado.
\end{solution}

\begin{solution}{exo:b2:ligand-field:6}
No \ce{[CoCl4]^2-} tetraédrico, o desdobramento é cerca de $4/9$ de um
desdobramento octaédrico, e o cloreto é um ligante mais fraco que a água: a
banda d--d se desloca para comprimentos de onda maiores (do laranja ao
vermelho), e o complexo parece azul. Os complexos tetraédricos também absorvem
mais fortemente, daí a cor intensa.
\end{solution}

\begin{solution}{exo:b2:ligand-field:7}
\ce{Ni^2+} é $d^8$. Quadrado plano: os quatro orbitais inferiores contêm os oito
elétrons em pares, e $\mathrm d_{x^2-y^2}$ fica vazio: \omterm{def:b2:diatomic-mos:magnetism}{diamagnético}. Tetraédrico:
$e^4t_2^4$, dois elétrons desemparelhados no conjunto $t_2$.
\end{solution}

\begin{solution}{exo:b2:ligand-field:8}
$\tilde\nu = 1/(500 \times 10^{-7}\,\text{cm}) = \qty{20000}{cm^{-1}}$; $\Delta_o =
0.1196/(500 \times 10^{-9}) = \qty{2.39e5}{J/mol}$, \qty{239}{kJ/mol}.
\end{solution}

\begin{solution}{exo:b2:ligand-field:9}
\ce{Zn^2+} é $d^{10}$ e \ce{Sc^3+} $d^0$: nenhuma \omterm{def:b2:ligand-field:d-d}{transição d--d} é possível (nenhum
nível d vazio para \ce{Zn^2+}, nenhum elétron d para \ce{Sc^3+}); portanto, não há
absorção no visível.
\end{solution}

\begin{solution}{exo:b2:ligand-field:10}
Co(III) $d^6$ + seis pares não ligantes da amônia: 18 elétrons. Doze preenchem os
seis \omterm{def:b2:diatomic-mos:bonding}{orbitais ligantes}, seis preenchem o $t_{2g}$ não ligante (a amônia dá um
campo forte o bastante para o \omterm{def:b2:ligand-field:spin}{spin baixo} com o cobalto(III)); $e_g^*$ está vazio.
Todos os elétrons estão emparelhados: \omterm{def:b2:diatomic-mos:magnetism}{diamagnético}.
\end{solution}

\begin{solution}{exo:b2:ligand-field:11}
O \ce{CO} é um bom doador $\sigma$ e, sobretudo, um aceptor $\pi$: seus orbitais
$\pi^*$ vazios abaixam o nível $t_{2g}$ por \omterm{def:b2:ligand-field:pi-ligands}{retrodoação}, aumentando $\Delta_o$. O
\ce{F-} é um doador $\pi$, que eleva $t_{2g}$ e diminui $\Delta_o$. O modelo de
cargas pontuais só enxerga a carga e erra a ordem.
\end{solution}

\begin{solution}{exo:b2:ligand-field:12}
Sem campo, as entalpias de hidratação seguiriam uma linha regular (os íons
encolhem ao longo da série). A água dá aquacomplexos de \omterm{def:b2:ligand-field:spin}{spin alto} cuja EECC é
nula para $d^0$, $d^5$, $d^{10}$ e máxima para $d^3$ e $d^8$: a estabilização
adicional acrescenta duas corcovas, com máximos perto de \ce{V^2+} e \ce{Ni^2+},
acima da reta que passa por \ce{Ca^2+}, \ce{Mn^2+}, \ce{Zn^2+}.
\end{solution}

\begin{solution}{pb:b2:ligand-field:1}
\textbf{1.} $d^3$.
\textbf{2.} $t_{2g}^3$, um elétron em cada um dos três orbitais, com spins paralelos.
\textbf{3.} Não: três elétrons cabem em $t_{2g}$ sem emparelhamento.
\textbf{4.} $-1.2\Delta_o$.
\textbf{5.} Três.
\textbf{6.} A EECC é grande e os orbitais $e_g^*$, que apontam para os ligantes,
estão vazios: retirar ou acrescentar um ligante custa boa parte dessa
estabilização.
\textbf{7.} Seis íons óxido nos vértices de um octaedro ligeiramente distorcido.
\textbf{8.} $1/(556 \times 10^{-7}) = \qty{17990}{cm^{-1}}$, cerca de \qty{1.80e4}{cm^{-1}}.
\textbf{9.} $0.1196/(556 \times 10^{-9}) = \qty{2.15e5}{J/mol}$: \qty{215}{kJ/mol}.
\textbf{10.} Amarelo-esverdeado (556 nm) e violeta (405 nm).
\textbf{11.} Vermelho, com um pouco de azul: o vermelho intenso do rubi.
\textbf{12.} Um íon $d^3$ tem mais de uma configuração excitada com um elétron em
$e_g$; a repulsão entre elétrons lhes dá energias diferentes, donde várias
bandas (tratadas no volume do 3.º ano).
\textbf{13.} $0.1196/(620 \times 10^{-9}) = \qty{1.93e5}{J/mol}$, \qty{193}{kJ/mol}
(\qty{16130}{cm^{-1}}).
\textbf{14.} O rubi: seu desdobramento é maior.
\textbf{15.} Para o vermelho-alaranjado: o vermelho passa a ser absorvido, e o
verde (com um pouco de azul) é transmitido.
\textbf{16.} No berilo, os sítios do cromo são um pouco maiores e os íons óxido
ficam mais afastados, e o entorno é diferente: um campo mais fraco.
\textbf{17.} Não: um íon $d^3$ conserva três elétrons desemparelhados em
qualquer campo octaédrico.
\textbf{18.} $d^9$.
\textbf{19.} O aquacomplexo absorve no vermelho e no laranja por uma \omterm{def:b2:ligand-field:d-d}{transição
d--d}: a luz transmitida é azul.
\textbf{20.} Ele fica branco: sem ligantes água em torno do cobre (o sulfato
toma seu lugar), o campo é mais fraco e a banda passa para o infravermelho.
\textbf{21.} A água se liga de novo ao cobre e o aquacomplexo azul se forma.
\textbf{22.} Em comprimento de onda menor: a amônia, acima da água na série, dá um
desdobramento maior (o azul-violeta intenso do complexo amoniacal).
\textbf{23.} Não: $t_{2g}^6e_g^3$ é a única configuração.
\textbf{24.} Sim: um elétron desemparelhado.
\textbf{25.} $\boldsymbol{\Delta_o \approx \qty{1.8e4}{cm^{-1}}}$, cerca de
$\boldsymbol{\qty{215}{kJ/mol}}$.
\end{solution}
